\scanpage{140}
\mspage{85}
% archive sheet number: 139

\[
\mathrm{Hot}^{\deleted{\unsure{\mathrm{lc}}}}_W(S)\longrightarrow\mathrm{Hot}^{\unsure{\mathrm{lc}|\mathrm{le}}}_W(S),
\]

mais pas de foncteur en sens inverse. Mais si, mais si, \deleted{\unsure{je vois poindre l'aspect géom. de ma construction}, car les chgts de base (lisses disons) qui m'avaient servi pour définir \deleted{\unsure{l'équ}} la $W$-équivalence fine sur $S$, \unsure{me donnent un certain chgt}} en travaillant avec la catégorie similaire \deleted{\unsure{de}} à $\underline{\mathrm{Fibhot}}\,S$, ça doit \unsure{garder|prendre} un sens, tout au moins quand on a une bonne notion de ``fibration''. Je vois un petit travail de dégrossissage conceptuel à faire, avec peut-être une formule d'adjonction à la \unsure{``clef|clef}, mais je ne vais pas m'y lancer à présent.

$\backslash$

\margin{\unsure{V.16|Ve.16|VI.16}}

Je vais quand même essayer de faire ce travail ``de dégrossissage'', car le diagramme des 15 catégories a l'air un peu touffu, finalement! Je me donne

\[
(\mathcal M,W)\qquad W\subset\mathrm{Fl}\,\mathcal M,
\]

un $S\in\mathrm{Ob}\,\mathcal M$, d'où $\mathcal M_{/S}$, avec $W_S\subset\mathrm{Fl}\,(\mathcal M_{/S})$ induit par $W$. On va poser

\[
\mathrm{Ho}^{\mathrm{lc}}(S)=(\mathcal M_{/S})W_S^{-1},
\]

ça jouera le rôle \unsure{du type d'homotopie relativement constant|des types d'homotopie relativement constants}. Je suppose d'autre part que l'on a distingué une classe $\unsure{\mathcal F|F}\subset\mathrm{Fl}\,(\mathcal M)$ de flèches, jouant le rôle de ``fibrations'' (ce seront les morphismes Hot-fibrants de $\mathrm{Cat}$, donnant lieu à la s.catégorie $\underline{\mathrm{Fibhot}}\,S$ de $\mathrm{Cat}_{/S}$). Je prévois qu'il faudra
